Una red de transporte sanitario de la Región Metropolitana debe decidir cada semana cuántos vehículos contratar, en qué bases ubicarlos y qué hubs habilitar, y cada día cómo atender entre 455 y 731 solicitudes semanales de muestras clínicas, fármacos refrigerados, insumos y equipos. Este informe propone una matheurística que coordina ambas escalas: heurísticas GRASP y VND generan rutas diarias, un modelo de particionamiento de conjuntos resuelto con Gurobi elige entre ellas la flota, las bases, los hubs y la forma de atender cada solicitud, y una generación de columnas con los precios duales del modelo agrega rutas que pueden bajar el costo. Se compara contra un caso base que representa la planificación por reglas.

Esta versión corrige el Informe~2 anterior: se corrigió un error del código que inutilizaba muchas rutas con solicitudes divididas, se implementó la generación de columnas con duales que antes solo se describía y se fortaleció el caso base, que ahora externaliza 10 solicitudes por semana en vez de 33,6.

En las 40 semanas BIG, la matheurística es más barata que el caso base en todas, con un ahorro promedio de 9,9\,\% (entre 8,0 y 12,9\,\%), equivalente a \$3,9 millones por semana. Casi todo el ahorro viene de rutear mejor: con prácticamente la misma flota recorre un 16\,\% menos de kilómetros y externaliza un 28\,\% menos. Todas las soluciones cumplen todas las restricciones según un validador independiente. Las solicitudes difíciles tienen un perfil claro: insumos y equipos grandes con ventanas de una hora, que necesitan furgones justo cuando los 18 disponibles están ocupados, lo que ocurre en 26 de las 40 semanas. En 22 corridas de sensibilidad (cambios de demanda, flota, congestión, jornada, costos e infraestructura) el ahorro se mantuvo entre 7,3 y 10,2\,\%; el modelo mantiene las nueve bases aun con su costo duplicado, cerrar Talagante no subió el costo y los hubs casi no aportan.

Las principales limitaciones son que la cota obtenida vale solo para las rutas generadas, que los tiempos de viaje suponen flujo libre y que la matheurística acepta más atrasos que el caso base porque su penalización es baja. Los pasos siguientes son medir el método en las instancias SMALL, extender la sensibilidad a más semanas y evaluar más capacidad en furgones y el cierre de la base de Talagante.
